\documentclass{article}
\usepackage[margin=0.25in]{geometry}
\usepackage[expansion=false]{microtype}
\usepackage{hyperref}
\usepackage{amsmath}
\usepackage{amssymb}
\usepackage{booktabs}
\usepackage{enumitem}
\setlist{nosep}

\title{Related Work Notes}
\author{Cameron Crow \and DrTrayaurus}
\date{}

\begin{document}
\maketitle

%% ============================================================
\section{Market Microstructure \& Price Instability}
%% ============================================================


%% ============================================================
\section{Information Asymmetry \& Lemon Markets}
%% ============================================================


%% ============================================================
\section{LLM Agents \& Multi-Agent Economics}
%% ============================================================


%% ============================================================
\section{Sybil Attacks \& Reputation Systems}
%% ============================================================


%% ============================================================
\section{Agent Benchmarks \& Evaluation}
%% ============================================================


%% ============================================================
\section{Vending Bench}
%% ============================================================

\textbf{Citation:} Andon Labs. ``Vending-Bench: A Benchmark for Long-Term Coherence of Autonomous Agents.'' arXiv:2502.15840, 2025.

\medskip
\noindent\textbf{Overview.}
VendingBench tasks a single LLM agent with running a vending machine business over a simulated multi-month horizon, managing pricing, inventory, and supplier negotiations. The primary metric is long-term cognitive coherence, measured by the agent's final bank balance. An informal extension called \emph{Arena} places multiple agents at the same location and introduces emergent price competition. VendingBench is the closest existing benchmark to AI Bazaar's competitive market setting, but it does not treat market failure or alignment breakdown as the research object---competitive dynamics are observed qualitatively in Arena but never benchmarked or analyzed.

\subsection*{Consumer Demand}

Demand is modeled as a \textbf{price-elastic daily oracle with stochastic noise}; there are no simulated consumer agents. Each product is assigned three GPT-4o-generated cached parameters---base sales, reference price, and price elasticity---at the start of the simulation. Daily sales volume is computed by scaling base sales by a price-deviation factor (using the elasticity coefficient), then multiplied by four stochastic modifiers: day-of-week, monthly/seasonal, weather impact, and a product-variety penalty (excess SKU count reduces demand, capped at a 50\% reduction). Gaussian noise is added before the result is rounded and clamped to available inventory. There is no closed-form demand equation published in the paper, no consumer budget constraint, and no agent-level buyer simulation; purchasing is entirely mediated through this aggregate formula.

\subsection*{Overhead Costs}

Overhead is minimal and uniform: a \textbf{flat \$2 per day} location/rental fee, the only recurring fixed cost beyond inventory procurement. Agents begin with \$500 in starting capital; bankruptcy is triggered if the fee goes unpaid for 10 consecutive days. VendingBench 2 adds a meta-cost: agents are billed for their own LLM output tokens at \$100 per million tokens per week, designed to penalize verbose reasoning. There is no variable overhead scaling with volume, no depreciation, and no tiered cost structure.

\subsection*{Firm--Firm Communication}

In the base single-agent setting there is \textbf{no competitor visibility or inter-firm communication}; the agent operates as a monopolist interacting only with simulated suppliers via email. The Arena extension introduces full transparency and direct coordination: agents can call \texttt{get\_machine\_inventory} with a competitor's email to retrieve slot-by-slot price and stock data, communicate via a shared email system, and transfer money between accounts. Documented emergent behaviors include explicit cartel proposals (``Should we agree on a price floor of \$2.00?''), exclusive supplier access negotiations, and side-payment offers to deter a competitor from restocking a category. This makes Arena substantially more communication-rich than AI Bazaar, where firms receive only market-level aggregates.

\subsection*{Market Knowledge Given to Firms}

Each morning in the base setting, the agent receives: (1) per-item \textbf{sales counts} from the previous day, (2) current \textbf{inventory state} via tool call, (3) \textbf{cash balance} via tool, and (4) \textbf{supplier emails}. Agents do \emph{not} receive competitor prices, market-level demand aggregates, or consumer behavior statistics---all market inference must be built from the agent's own sales history. Agents additionally have internet search access (via Perplexity) to research real-world wholesale pricing. The system prompt evaluates the agent solely on its final bank balance. In Arena, information expands substantially: agents can query competitor machines on demand, giving full real-time competitor pricing and inventory access.

\subsection*{Key Contrasts with AI Bazaar}

\begin{center}
\begin{tabular}{lll}
\toprule
Dimension & VendingBench & AI Bazaar \\
\midrule
Primary metric & Long-term coherence / final balance & Market stability / alignment failure \\
Demand model & Price-elastic daily oracle with noise & CES consumers, Poisson demand \\
Fixed overhead & \$2/day flat fee & Configurable (daily scaled) \\
Competitor visibility & None (base); full query (Arena) & Market aggregates only \\
Inter-firm communication & Email + payment transfers (Arena) & No direct channel \\
Failure modes studied & Meltdown, forgetting, supplier scams & THE\_CRASH, LEMON\_MARKET \\
Multi-agent framing & Informal Arena (non-benchmark) & Controlled benchmark \\
\bottomrule
\end{tabular}
\end{center}

\noindent VendingBench does not model or study systemic market failure, price war dynamics, or alignment failures as research objects. AI Bazaar is therefore the first work to treat these competitive dynamics as the primary phenomenon under study.

\end{document}
